% !TeX program = xelatex
% Anleitung für Lehrende: OpenCode (Desktop-App) installieren und mit KI:connect verbinden
% Allgemein gehalten, weil sie für mehrere Workshops des ZWEK dient.
% Quellen (10.11.2026): opencode.ai/download (Desktop-App für macOS, Windows, GNU/Linux),
% Durchlauf mit OpenCode Desktop 2.2.7 auf einem MacBook (macOS 26.3) in einem Testprofil:
% Projekt hinzufügen, Ordner im macOS-Dialog anlegen und öffnen, Starterordner liest
% opencode.json (KI:connect / GPT-OSS 120B vorgewählt), /connect mit Platzhalter-Schlüssel.
% Rote Rahmen in den Bildern mit ImageMagick.
% Nicht geprüft: Windows und GNU/Linux (Dialoge sehen anders aus), Antwort auf die Testfrage
% mit echtem Schlüssel in der Desktop-App, SmartScreen- und Gatekeeper-Hinweise (stehen als „falls“).
% Gestaltung nach den Moodle-Hilfen des ZWEK: Rot für Schrittnummern und Achtung,
% Grau für Checklisten und Tipps, eckige Formen, Icons als PNG.
\documentclass[a4paper,11pt]{article}

\usepackage{fontspec}
\IfFontExistsTF{Arial}
  {\setmainfont{Arial}\setsansfont{Arial}}
  {\setmainfont{Liberation Sans}\setsansfont{Liberation Sans}}
\setmonofont{DejaVu Sans Mono}[Scale=0.88]
\newcommand{\versal}[1]{{\addfontfeature{LetterSpace=6.0}\MakeUppercase{#1}}}

\usepackage[ngerman]{babel}
\usepackage{xcolor}
\definecolor{hsdrot}{HTML}{E60028}
\definecolor{hsdgrau}{HTML}{58585A}
\definecolor{hsddunkel}{HTML}{333333}
\definecolor{hsdbalken}{HTML}{303031}
\definecolor{hsdhell}{HTML}{F2F2F2}
\definecolor{hsdmittel}{HTML}{D9D9D9}
\definecolor{hsdlinie}{HTML}{BEBEBE}

\usepackage[a4paper,left=2.8cm,right=2.8cm,top=2.2cm,bottom=2.6cm,footskip=1.1cm,headheight=14pt]{geometry}
\usepackage{graphicx}
\usepackage{tikz}
\usepackage[most]{tcolorbox}
\usepackage{fancyvrb}
\usepackage{enumitem}
\usepackage{ragged2e}
\usepackage{needspace}
\usepackage{fancyhdr}
\usepackage{lastpage}
\usepackage{ccicons}
\usepackage{microtype}
\usepackage[hidelinks]{hyperref}
\usepackage{hyperxmp}
\urlstyle{same}
\hypersetup{pdftitle={OpenCode einrichten: Anleitung für Lehrende},
  pdfauthor={Christian-Maximilian Steier (ZWEK\xmpcomma{} Hochschule Düsseldorf)},
  pdfsubject={OpenCode Desktop installieren und mit KI:connect.nrw verbinden},
  pdfkeywords={OpenCode, KI:connect.nrw, Anleitung, Lehrende, ZWEK, Hochschule Düsseldorf, CC BY 4.0},
  pdfcopyright={© 2026 Christian-Maximilian Steier, ZWEK, Hochschule Düsseldorf. Lizenz CC BY 4.0, https://creativecommons.org/licenses/by/4.0/},
  pdflicenseurl={https://creativecommons.org/licenses/by/4.0/},
  pdflang={de-DE},pdfdisplaydoctitle=true}
% hyperxmp unter XeTeX schreibt Autor und Schlagworte nur ins XMP, nicht ins Info-Verzeichnis.
% Klartext: Christian-Maximilian Steier (ZWEK, Hochschule Düsseldorf) | OpenCode, KI:connect.nrw, Anleitung, Lehrende, ZWEK, Hochschule Düsseldorf, CC BY 4.0
\AtBeginDocument{\special{pdf:docinfo << /Author <FEFF00430068007200690073007400690061006E002D004D006100780069006D0069006C00690061006E002000530074006500690065007200200028005A00570045004B002C00200048006F006300680073006300680075006C00650020004400FC007300730065006C0064006F007200660029> /Keywords <FEFF004F00700065006E0043006F00640065002C0020004B0049003A0063006F006E006E006500630074002E006E00720077002C00200041006E006C0065006900740075006E0067002C0020004C0065006800720065006E00640065002C0020005A00570045004B002C00200048006F006300680073006300680075006C00650020004400FC007300730065006C0064006F00720066002C00200043004300200042005900200034002E0030> >>}}
\graphicspath{{icons/}{Bilder/}}

\renewcommand{\normalsize}{\fontsize{11}{13.6}\selectfont}
\normalsize
\RaggedRight
\hyphenpenalty=10000 \exhyphenpenalty=10000
\emergencystretch=3em
\setlength{\parindent}{0pt}
\setlength{\parskip}{6pt}
\raggedbottom
\clubpenalty=10000 \widowpenalty=10000
\setlist{itemsep=3pt,topsep=3pt,leftmargin=1.4em}

\newcommand{\dokstand}{11.11.2026}
\newcommand{\zwekmail}{\href{mailto:zwek@hs-duesseldorf.de}{zwek@hs-duesseldorf.de}}
\pagestyle{fancy}
\fancyhf{}
\renewcommand{\headrulewidth}{0pt}
\renewcommand{\footrulewidth}{0.6pt}
\fancyhead[L]{\footnotesize\bfseries\color{hsdgrau}OpenCode einrichten\quad\textcolor{hsdrot}{\textperiodcentered}\quad Anleitung für Lehrende}
\fancyhead[R]{\raisebox{-2pt}{\textcolor{hsdrot}{\rule{1.4cm}{2.2pt}}}}
\fancyfoot[L]{\footnotesize\color{hsdgrau}\textbf{ZWEK \textbar\ Hochschule Düsseldorf}\\
  Stand \dokstand\ \textperiodcentered\ gültig für OpenCode Desktop~2.2 \textperiodcentered\ Pflege: ZWEK, Hochschul- und Mediendidaktik}
\fancyfoot[R]{\footnotesize\color{hsdgrau}\hyperref[lizenz]{\raisebox{-0.1em}{\ccby}\;CC~BY~4.0}\quad\textbf{Seite \thepage\ von \pageref*{LastPage}}}
\fancypagestyle{erste}{\fancyhead{}}

\newcommand{\icon}[2][0.95em]{\raisebox{-0.14em}{\includegraphics[height=#1]{#2.png}}}

% Schrittkopf: großer roter Kreis mit weißer Ziffer, darüber grau „Schritt n von 5 · Dauer“
\newcommand{\schritt}[3]{%
  \par\needspace{9\baselineskip}\vspace{16pt}%
  \noindent\begin{minipage}[c]{40pt}\begin{tikzpicture}
    \node[circle,fill=hsdrot,text=white,minimum size=30pt,inner sep=0pt,
      font=\fontsize{17}{17}\selectfont\bfseries]{#1};
  \end{tikzpicture}\end{minipage}%
  \begin{minipage}[c]{\dimexpr\linewidth-40pt\relax}
    {\footnotesize\bfseries\color{hsdgrau}\versal{Schritt #1 von 5}\enspace\textperiodcentered\enspace\mdseries #3}\par\vspace{1pt}
    {\fontsize{17}{20}\selectfont\bfseries\color{hsddunkel}#2\par}
  \end{minipage}\par
  \nobreak\vspace{4pt}{\color{hsdmittel}\rule{\linewidth}{0.6pt}}\par\nobreak\vspace{2pt}\nobreak}

% kleine Kreise für Teilschritte in Karten
\newcommand{\nr}[1]{\tikz[baseline=(n.base)]\node[circle,fill=hsdrot,text=white,inner sep=0pt,
  minimum size=13pt,font=\scriptsize\bfseries](n){#1};}
\newenvironment{teilschritte}
  {\begin{list}{}{\setlength{\leftmargin}{20pt}\setlength{\labelwidth}{14pt}
   \setlength{\labelsep}{6pt}\setlength{\itemsep}{4pt}\setlength{\topsep}{4pt}}}
  {\end{list}}
\newcommand{\ts}[1]{\item[\nr{#1}]}

% Tasten und Oberflächenbegriffe: dünner grauer Rahmen
\newcommand{\taste}[1]{\tcbox[on line,boxsep=0.5pt,left=4pt,right=4pt,top=1.5pt,bottom=1.5pt,
  colback=white,colframe=hsdlinie,boxrule=0.6pt,arc=0pt]{\footnotesize\bfseries #1}}
\newcommand{\datei}[1]{\texttt{#1}}
\newcommand{\weiter}{\,\tikz[baseline=-0.6ex]\fill[hsdrot](0,-0.09)--(0.13,0)--(0,0.09)--cycle;\,}

% Befehlsblock: dunkle Fläche, weiße Schrift, Kennzeile „Eingeben“
\fvset{formatcom=\color{white},fontsize=\fontsize{12}{15}\selectfont}
\newtcolorbox{befehl}[1][Eingeben]{enhanced,breakable=false,arc=0pt,boxrule=0pt,
  colback=hsddunkel,colframe=hsddunkel,left=10pt,right=10pt,top=5pt,bottom=7pt,
  before skip=6pt,after skip=8pt,
  borderline west={4pt}{0pt}{hsdrot},
  before upper={\setlength{\parskip}{0pt}\setlength{\topsep}{0pt}\setlength{\partopsep}{0pt}%
    {\scriptsize\bfseries\color{hsdmittel}\icon[0.85em]{terminal-weiss}\;\versal{#1}}\par\vspace{-5pt}},
  after upper={\vspace{-4pt}}}

% Kästen nach den Moodle-Hilfen: grau = Zusatz, roter Rand = Achtung
\tcbset{hsdkasten/.style={enhanced,arc=0pt,boxrule=0pt,colback=hsdhell,colframe=hsdhell,
  left=13pt,right=12pt,top=9pt,bottom=9pt,before skip=10pt,after skip=10pt}}
\newtcolorbox{checkliste}[1]{hsdkasten,borderline west={4pt}{0pt}{hsdlinie},
  before upper={{\footnotesize\bfseries\color{hsddunkel}\icon{check-circle-grau}\;\versal{#1}}\par\vspace{5pt}}}
\newtcolorbox{tipp}{hsdkasten,borderline west={4pt}{0pt}{hsdlinie},
  before upper={{\footnotesize\bfseries\color{hsddunkel}\icon{lightbulb-grau}\;\versal{Tipp}}\par\vspace{3pt}}}
\newtcolorbox{achtung}[1]{hsdkasten,borderline west={4pt}{0pt}{hsdrot},
  before upper={{\footnotesize\bfseries\color{hsdrot}\icon{exclamation-triangle-rot}\;\versal{Achtung}}\par\vspace{3pt}{\bfseries #1}\par\vspace{2pt}}}

\newcommand{\quelle}{https://github.com/HSD-ZWEK/Resources/tree/main/KI/OpenCode}
\newcommand{\starterzip}{https://github.com/HSD-ZWEK/Resources/raw/main/KI/OpenCode/starterordner-ki-wissen.zip}
\newtcolorbox{lizenz}{hsdkasten,borderline west={4pt}{0pt}{hsdlinie},
  before upper={{\footnotesize\bfseries\color{hsddunkel}\versal{Weitergeben und anpassen erlaubt}}\par\vspace{5pt}}}
\newcommand{\kaestchen}{\tikz[baseline=-0.8ex]\draw[hsdgrau,line width=0.8pt](0,-0.15)rectangle(0.28,0.13);}
\newcommand{\haken}{\tikz[baseline=-0.8ex]{\draw[hsdgrau,line width=0.8pt](0,-0.15)rectangle(0.28,0.13);
  \draw[hsdrot,line width=1.3pt](0.05,0)--(0.12,-0.08)--(0.25,0.1);}}
\newenvironment{haekchen}[1]
  {\begin{list}{}{\setlength{\leftmargin}{20pt}\setlength{\labelwidth}{14pt}
   \setlength{\labelsep}{6pt}\setlength{\itemsep}{4pt}\setlength{\topsep}{0pt}}
   \let\punkt\relax\def\punkt{\item[#1]}}
  {\end{list}}

% Karte mit dunklem Titelbalken und roter Akzentlinie
\newtcolorbox{karte}[2]{enhanced,arc=0pt,boxrule=0.6pt,colframe=hsdmittel,colback=white,
  colbacktitle=hsdbalken,coltitle=white,fonttitle=\bfseries,
  title={\icon[1.05em]{#1}\;\;#2},toptitle=4pt,bottomtitle=4pt,left=10pt,right=10pt,
  top=7pt,bottom=8pt,before skip=8pt,after skip=8pt,
  overlay={\draw[hsdrot,line width=2pt]([yshift=1pt]title.south west)--([yshift=1pt]title.south east);}}

% Bildschirmfoto: 1 px grauer Rahmen (HSD-Regel), Unterschrift darunter
\newcommand{\bild}[3][\linewidth]{\par\vspace{4pt}\noindent
  \begin{minipage}{\linewidth}
  \fcolorbox{hsdlinie}{white}{\includegraphics[width=\dimexpr#1-2\fboxsep-2\fboxrule\relax]{#2}}\par\vspace{2pt}
  {\footnotesize\color{hsdgrau}#3}
  \end{minipage}\par\vspace{6pt}}
\setlength{\fboxsep}{0pt}

% Teilschritt und Bild bleiben auf einer Seite
\newcommand{\tsbild}[5]{\par\noindent\begin{minipage}{\linewidth}
  \begin{teilschritte}\ts{#1} #2\end{teilschritte}\bild[#3]{#4}{#5}\end{minipage}\par\vspace{2pt}}

% Problem-Lösung-Karte
\newtcolorbox{problem}[1]{enhanced,arc=0pt,boxrule=0pt,colback=hsdhell,colframe=hsdhell,
  borderline west={4pt}{0pt}{hsdlinie},left=11pt,right=10pt,top=7pt,bottom=8pt,
  before upper={{\bfseries\color{hsddunkel}#1}\par\vspace{3pt}%
    {\footnotesize\bfseries\color{hsdgrau}\versal{Lösung}}\par\vspace{1pt}}}

\begin{document}
\thispagestyle{erste}

% ------------------------------------------------------------------ Kopf
\noindent\includegraphics[width=3.4cm]{hsd-logo.png}\hfill
\raisebox{0.9cm}{\footnotesize\bfseries\color{hsdgrau}\versal{Anleitung für Lehrende}}
\par\vspace{4pt}
{\color{hsdrot}\rule{\linewidth}{2.2pt}}\par\vspace{10pt}
{\fontsize{26}{30}\selectfont\bfseries OpenCode einrichten\par}
\vspace{4pt}
{\large\color{hsddunkel}und mit KI:connect.nrw verbinden. Etwa 15 Minuten.\par}
\vspace{10pt}

\textbf{OpenCode ist ein Programm, das auf dem eigenen Rechner installiert ist.}
Das Verfahren erfolgt unter Zuhilfenahme eines Ordners, der eine Vielzahl an
Dokumenten enthält. Die gestellte Frage wird in Verbindung mit den
entsprechenden Stellen an KI:connect.nrw übermittelt. Der Antwortentwurf wird
als Datei in den betreffenden Ordner zurückgeschrieben. Die Verarbeitung
erfolgt auf Hochleistungsrechnern an nordrhein-westfälischen Hochschulen.

\begin{checkliste}{Das brauchen Sie}
\begin{haekchen}{\kaestchen}
\punkt Einen Rechner mit \textbf{Windows, macOS oder GNU/Linux}, auf dem Sie Programme installieren dürfen
\punkt Rund \textbf{1~GB} freien Speicher und Internet
\punkt Den \textbf{Starterordner} vom ZWEK: \href{\starterzip}{starterordner-ki-wissen.zip}
\punkt Ein Konto bei \textbf{KI:connect.nrw}. Sie melden sich dort mit Ihrer HSD-Kennung an. Den Schlüssel legen Sie in Schritt~4 selbst an.
\end{haekchen}
\end{checkliste}

% ------------------------------------------------------------------ Wegleiste
\vspace{4pt}
{\footnotesize\bfseries\color{hsddunkel}\versal{Ihr Weg in fünf Schritten}}\par\vspace{6pt}
\noindent\begin{tikzpicture}[x=\linewidth/5]
  \draw[hsdmittel,line width=2pt](0.5,0)--(4.5,0);
  \foreach \i/\t in {1/Installieren,2/Ordner anlegen,3/Starterdateien hineinlegen,4/Schlüssel anlegen,5/Testfrage}{
    \node[circle,fill=hsdrot,text=white,minimum size=22pt,inner sep=0pt,
      font=\normalsize\bfseries] at (\i-0.5,0) {\i};
    \node[below,text width=0.19\linewidth,align=center,font=\footnotesize\bfseries,
      text=hsddunkel] at (\i-0.5,-0.5) {\t};}
\end{tikzpicture}

\begin{tipp}
Laden Sie OpenCode nur von \textbf{opencode.ai/download}. Auf anderen Seiten,
auch auf GitHub, finden sich ältere Fassungen.
\end{tipp}

% ------------------------------------------------------------------ Schritt 1
\schritt{1}{OpenCode installieren}{etwa 5 Minuten}

Öffnen Sie im Browser \href{https://opencode.ai/download}{\textbf{opencode.ai/download}}
und laden Sie im Bereich \textbf{Desktop} die Datei für Ihr System.

\begin{karte}{apple-weiss}{macOS}
\begin{teilschritte}
\ts{1} \textbf{macOS (Apple Silicon)} für Macs ab 2020 mit M-Chip, sonst
\textbf{macOS (Intel)}. Sie bekommen eine \datei{.dmg}-Datei.
\ts{2} Die Datei doppelklicken und \textbf{OpenCode} auf den Ordner
\textbf{Programme} ziehen.
\ts{3} OpenCode aus dem Ordner Programme starten.
\end{teilschritte}
\end{karte}

\begin{tcbraster}[raster columns=2,raster equal height,raster column skip=10pt,
  raster before skip=2pt,raster after skip=6pt]
\begin{karte}{windows-weiss}{Windows}
\begin{teilschritte}
\ts{1} \textbf{Windows (x64)}. Nur Geräte mit Snapdragon-Chip brauchen
\textbf{Windows (ARM64)}.
\ts{2} Die \datei{.exe} doppelklicken und der Installation folgen.
\end{teilschritte}
\end{karte}
\begin{karte}{linux-weiss}{GNU/Linux}
\begin{teilschritte}
\ts{1} Ubuntu und Debian: \datei{.deb}. Fedora: \datei{.rpm}.
\ts{2} Die Datei doppelklicken, die Softwareverwaltung installiert sie.
\end{teilschritte}
\end{karte}
\end{tcbraster}

Beim ersten Start sehen Sie die Startseite mit der Überschrift
\textbf{Projekte}. Darunter steht der Name Ihres Rechners.

\begin{tipp}
Falls macOS meldet, die App könne nicht geöffnet werden: Systemeinstellungen
\weiter{} Datenschutz \& Sicherheit, nach unten scrollen, \textbf{Dennoch öffnen}.
Falls Windows einen blauen SmartScreen-Hinweis zeigt: \textbf{Weitere
Informationen} \weiter{} \textbf{Trotzdem ausführen}.
\end{tipp}

% ------------------------------------------------------------------ Schritt 2
\schritt{2}{Den Ordner anlegen und öffnen}{etwa 2 Minuten}

OpenCode arbeitet immer in einem Ordner, dem \textbf{Projekt}. Den Ordner legen
Sie direkt in OpenCode an.

\tsbild{1}{Mit der Maus über den Namen Ihres Rechners fahren. Rechts erscheint ein
Ordnersymbol mit Pluszeichen, \textbf{Projekt hinzufügen}. Darauf klicken.}{0.45\linewidth}{01-projekt-hinzufuegen.png}{Hier heißt der Rechner „Arbeitslaptop“. Bei Ihnen steht dort der Name Ihres Rechners.}

\tsbild{2}{Ein Fenster zur Ordnerwahl öffnet sich. Links \textbf{Dokumente}
wählen, dann unten \textbf{Neuer Ordner}.}{0.66\linewidth}{02-neuer-ordner.png}{Der Ordnerdialog unter macOS. Unter Windows heißt die Schaltfläche ebenfalls „Neuer Ordner“.}

\tsbild{3}{Als Namen \datei{ki-wissen} eintippen und \textbf{Erstellen} klicken.}{0.66\linewidth}{03-name-erstellen.png}{Der neue Ordner liegt in Ihren Dokumenten.}

\tsbild{4}{Das Fenster steht jetzt im neuen Ordner \datei{ki-wissen}.
\textbf{Öffnen} klicken.}{0.58\linewidth}{04-oeffnen.png}{Oben in der Mitte steht der Name des Ordners, der geöffnet wird.}

\tsbild{5}{Unter Ihrem Rechnernamen steht jetzt \textbf{ki-wissen}. Rechts
erscheint \textbf{Neue Sitzung}. Damit geht es in Schritt~4 weiter.}{0.85\linewidth}{05-neue-sitzung.png}{Das Projekt ki-wissen ist angelegt.}

% ------------------------------------------------------------------ Schritt 3
\schritt{3}{Die Starterdateien hineinlegen}{etwa 3 Minuten}

Der Starterordner vom ZWEK enthält Musterfakten, Beispielanfragen und die
Einstellungen für KI:connect. Alles darin ist erfunden oder ein Muster.

\begin{teilschritte}
\ts{1} \href{\starterzip}{\datei{starterordner-ki-wissen.zip}} laden und doppelklicken
(Windows: rechte Maustaste \weiter{} \textbf{Alle extrahieren}).
\ts{2} Den Ordner \datei{ki-wissen} im Finder oder Explorer öffnen. Er liegt in
Ihren Dokumenten. In OpenCode geht es auch so: Maus über \textbf{ki-wissen},
drei Punkte \weiter{} \textbf{Im Finder anzeigen}.
\ts{3} Den entpackten Starterordner öffnen. Unter macOS
\taste{cmd}\,+\,\taste{Shift}\,+\,\taste{.} drücken, unter GNU/Linux
\taste{Strg}\,+\,\taste{H}. Dann erscheint auch der Ordner \datei{.opencode}.
Unter Windows ist er schon sichtbar.
\ts{4} Alles markieren mit \taste{cmd}\,+\,\taste{A} beziehungsweise
\taste{Strg}\,+\,\taste{A} und in den Ordner \datei{ki-wissen} in Ihren
Dokumenten ziehen.
\end{teilschritte}

\begin{tcbraster}[raster columns=2,raster equal height,raster column skip=10pt,
  raster before skip=4pt,raster after skip=4pt]
\begin{tcolorbox}[enhanced,arc=0pt,boxrule=0.6pt,colframe=hsdmittel,colback=white,
  left=10pt,right=10pt,top=7pt,bottom=7pt]
{\footnotesize\bfseries\color{hsddunkel}\versal{Danach liegt in ki-wissen}}\par\vspace{3pt}
\datei{.opencode} Arbeitsanweisungen\\
\datei{anfragen} Beispielanfragen\\
\datei{entwuerfe} hier landen die Entwürfe\\
\datei{wissen} Musterfakten und Vorlage\\
\datei{AGENTS.md}, \datei{opencode.json}
\end{tcolorbox}
\begin{tcolorbox}[enhanced,arc=0pt,boxrule=0.6pt,colframe=hsdmittel,colback=white,
  left=10pt,right=10pt,top=7pt,bottom=7pt]
{\footnotesize\bfseries\color{hsddunkel}\versal{Bitte unverändert lassen}}\par\vspace{3pt}
\datei{opencode.json} legt fest: OpenCode spricht nur mit KI:connect, führt
keine Befehle aus und schreibt nur innerhalb des Ordners.
\end{tcolorbox}
\end{tcbraster}

% ------------------------------------------------------------------ Schritt 4
\schritt{4}{Den Schlüssel für KI:connect anlegen und eintragen}{etwa 4 Minuten, einmalig}

Zuerst legen Sie im Browser Ihren persönlichen Schlüssel an.

\begin{teilschritte}
\ts{1} \href{https://chat.kiconnect.nrw}{\datei{chat.kiconnect.nrw}} öffnen und mit der HSD-Kennung anmelden.
\end{teilschritte}

\tsbild{2}{Unten links auf Ihren Namen klicken und \textbf{API-Schlüsselverwaltung}
wählen.}{0.5\linewidth}{11-kiconnect-menue.png}{Hier steht „Vorname Nachname“. Bei Ihnen steht dort Ihr Name.}

\tsbild{3}{Unter \textbf{Schlüsselname} einen Namen eintragen, etwa \datei{OpenCode}.
Daneben steht Ihre Nutzergruppe \textbf{Beschäftigte und Studierende der HSD}. Die
lassen Sie so. Dann \textbf{Schlüssel erstellen} klicken und den neuen Schlüssel
kopieren.}{0.75\linewidth}{12-kiconnect-schluessel.png}{Mit dem Mülleimer in der Liste löschen Sie einen Schlüssel wieder.}

Dann tragen Sie ihn in OpenCode ein.


\tsbild{4}{In OpenCode auf \textbf{Neue Sitzung} klicken (Ende von Schritt~2). Unten im
Eingabefeld steht jetzt \textbf{KI:connect / GPT-OSS 120B}. Das hat der
Starterordner eingestellt.}{\linewidth}{09-modell.png}{Steht hier etwas anderes, fehlen die Starterdateien im Ordner (Schritt~3).}

\tsbild{5}{In das Eingabefeld \datei{/connect} tippen und \taste{Enter} drücken.}{\linewidth}{06-connect.png}{Über dem Feld erscheint der passende Befehl.}

\tsbild{6}{Im Fenster \textbf{Integrations} oben in die Suche \datei{KI:connect}
tippen und den Eintrag \textbf{KI:connect} anklicken.}{0.62\linewidth}{07-ki-connect-waehlen.png}{Die Liste ist lang. Mit der Suche geht es schneller.}

\tsbild{7}{Den kopierten Schlüssel in das Feld einfügen und
\textbf{Weiter} klicken. Im nächsten Fenster ist \textbf{GPT-OSS 120B} schon
angehakt. Noch einmal \textbf{Weiter}.}{0.62\linewidth}{08-schluessel.png}{Unten rechts meldet OpenCode „KI:connect verbunden“.}

OpenCode legt den Schlüssel in seinem eigenen Speicher Ihres Benutzerkontos ab
und nicht im Ordner \datei{ki-wissen}.

\begin{achtung}{Ihr Schlüssel ist persönlich.}
Das Feld zeigt den Schlüssel im Klartext. Tragen Sie ihn nicht ein, während
jemand zusieht oder Ihr Bildschirm geteilt wird. Wer ihn hat, kann in Ihrem
Namen Anfragen an KI:connect stellen. Schreiben Sie ihn nie in eine Datei,
einen Chat oder eine Mail. Ist er versehentlich weitergegeben worden, löschen
Sie ihn in der API-Schlüsselverwaltung und legen einen neuen an.
\end{achtung}

% ------------------------------------------------------------------ Schritt 5
\schritt{5}{Testfrage stellen}{etwa 1 Minute}

Tippen Sie in das Eingabefeld und klicken Sie auf den Pfeil rechts:

\begin{befehl}[Eintippen]
\begin{Verbatim}
Wann hat das Prüfungsamt in der Vorlesungszeit Sprechzeit?
\end{Verbatim}
\end{befehl}
\bild{10-testfrage.png}{Der Pfeil schickt die Frage ab. \taste{Enter} geht auch.}

OpenCode liest die Datei \datei{wissen/pruefungsamt.md} und antwortet mit den
Sprechzeiten aus dem Muster.

\begin{checkliste}{Fertig, wenn \ldots}
\begin{haekchen}{\haken}
\punkt links unter Ihrem Rechnernamen das Projekt \textbf{ki-wissen} steht
\punkt unten im Eingabefeld \textbf{KI:connect / GPT-OSS 120B} steht
\punkt die Testfrage mit den \textbf{Sprechzeiten aus dem Muster} beantwortet wird
\end{haekchen}
\end{checkliste}

Beim nächsten Mal starten Sie OpenCode, klicken links auf \textbf{ki-wissen}
und dann auf \textbf{Neue Sitzung}. Der Schlüssel bleibt gespeichert.

% ------------------------------------------------------------------ Hilfe
\par\needspace{14\baselineskip}\vspace{16pt}
\noindent{\fontsize{17}{20}\selectfont\bfseries\color{hsddunkel}\icon[0.9em]{wrench-grau}\;\;Wenn etwas klemmt}\par
\nobreak\vspace{2pt}{\color{hsdmittel}\rule{\linewidth}{0.6pt}}\par\nobreak\vspace{2pt}

\begin{tcbraster}[raster columns=2,raster equal height=rows,raster column skip=10pt,
  raster row skip=8pt,raster before skip=6pt]
\begin{problem}{„Beim Starten des lokalen Servers ist ein Fehler aufgetreten“}
OpenCode beenden und neu starten. Läuft OpenCode zusätzlich im Terminal, dieses
zuerst beenden.
\end{problem}
\begin{problem}{Unten steht nicht „KI:connect / GPT-OSS 120B“}
Die Starterdateien fehlen im Ordner, oft der versteckte Ordner oder
\datei{opencode.json}. Schritt~3 wiederholen.
\end{problem}
\begin{problem}{In der Liste fehlt KI:connect}
Eine Sitzung in einem anderen Projekt ist offen. Links \textbf{ki-wissen}
wählen, \textbf{Neue Sitzung}, dann \datei{/connect}.
\end{problem}
\begin{problem}{„Schlüssel ungültig“ oder Fehler 401}
Schlüssel noch einmal eintragen (Schritt~4). Auf Leerzeichen am Anfang oder
Ende achten. Schlüssel laufen nach einiger Zeit ab. Dann in der
API-Schlüsselverwaltung einen neuen erstellen und eintragen.
\end{problem}
\begin{problem}{Die Antwort nennt keine Sprechzeiten}
Prüfen, ob \datei{wissen/pruefungsamt.md} im Ordner \datei{ki-wissen} liegt.
Dann eine neue Sitzung beginnen.
\end{problem}
\begin{problem}{Projekt im falschen Ordner angelegt}
Mit der Maus über das Projekt fahren, drei Punkte \weiter{} \textbf{Schließen}.
Dann Schritt~2 wiederholen.
\end{problem}
\end{tcbraster}

\begin{problem}{Nichts davon hilft}
Schicken Sie ein Bildschirmfoto der Meldung an \zwekmail.
Bitte ohne den Schlüssel.
\end{problem}

\begin{lizenz}\label{lizenz}
\begin{minipage}[t]{2.3cm}\vspace{0pt}{\fontsize{26}{26}\selectfont\color{hsddunkel}\ccby}\end{minipage}%
\begin{minipage}[t]{\dimexpr\linewidth-2.3cm}\vspace{0pt}\small
Diese Anleitung steht unter der Lizenz CC~BY~4.0. Sie dürfen sie kopieren,
weitergeben und für eigene Kurse anpassen, auch Bilder daraus. Nennen Sie dabei
die Quelle, am einfachsten mit diesem Satz:

\vspace{5pt}
\colorbox{white}{\parbox{\dimexpr\linewidth-2\fboxsep}{\footnotesize\RaggedRight
„OpenCode einrichten. Anleitung für Lehrende“ von Christian-Maximilian Steier,
ZWEK, Hochschule Düsseldorf, Lizenz CC~BY~4.0
(\url{https://creativecommons.org/licenses/by/4.0/deed.de}),
Quelle: \url{\quelle}}}

\vspace{5pt}\footnotesize\color{hsdgrau}
Haben Sie etwas geändert, ergänzen Sie „bearbeitet von …“. Die Logos der
Hochschule und von OpenCode sind von der Lizenz ausgenommen.
\end{minipage}
\end{lizenz}


\end{document}
